

\documentclass[../main.tex]{subfiles}

\begin{document}
\newpage

\begin{center}
\textbf{\LARGE Abstract}
\end{center}

PuckTrick is an open-source Python library developed at the University of Milano-Bicocca, designed to introduce controlled errors into tabular datasets intended for machine learning applications. The library enables the simulation of five types of data corruption: duplicates, inconsistent labels, missing values, noise, and outliers. Each error type can be configured through a declarative strategy defined in dictionary format, allowing flexible and reproducible data perturbation.

This work was conducted within the Spark\_Conversion branch and was firstly focused on extending PuckTrick to support Apache Spark (PySpark) as an alternative processing engine to Pandas. The objective of this extension is to enable the library to operate efficiently on large-scale datasets in a distributed environment, while preserving the same functionalies and ensuring functionally equivalent results across the two backends.

Following the implementation phase, the extended version of PuckTrick was evaluated on real-world data. The goal of this experimental phase was to assess how the synthetic noise generated by the library may influence (both positively and negatively) the performance of deep learning models trained on a cybersecurity dataset containing 16 million instances, which was intentionally corrupted using PuckTrick. This in hope of finding out what types of noises are most damaging for this kind of task and evetually if it possible to increase models perfomance by selectively injecting noise inside Cybersecurity domain's dataset.



\end{document}